\chapter{Evaluation Methodology}
This chapter specifies how the proposed metric will be evaluated empirically. It connects the research questions to the selected datasets, models, implementation variants, baselines, and statistical procedures used to test the stated hypotheses.

\section{Research Questions}
Defining research questions and hypothesis like "does FELICS behave in the way we expect it".

\section{Datasets}
Explain datasets and why they are chosen for evaluation: 1. BBQ, 2. MedQA, 3. NewGermanCredit

\section{Evaluated Models}
Explain which models and why I am testing them. Also talk briefly about constraints that led to the decision.


\section{Reference Implementation}
Explain things about the actual implementation, where to find it (Github) and that its essentially a fork from Mattons original. Also explain that we also run Mattons implementation as comparison with the same settings and where to find their implementation (also on GitHub).

\section{Baselines}
What am I comparing against: Once FELICS with correlation enabled ($max_in_group = 4$), FELICS with correlation disabled ($max_ib_group = 1$) and Mattons original.

\section{Evaluation Metrics}
How am I verifying my questions and hypothesis? E.g. conducting a t-test on results in R to check for significant changes, to practically prove backwards compatibility.

\chapter{Results}
The most pressing results in this chapter. Move detailed views and large tables into Appendix. Therefore keep this short, max 3 pages. Just mention observations, but keep explanations or discussions for later.

\chapter{Discussion}
This chapter interprets the empirical findings in relation to the research questions and the original design objectives. It also distinguishes limitations of the experimental setup from limitations inherent to the metric and uses both to motivate directions for future work.

\section{Interpretation}
Look based on the data if research questions can be answered and what the result is. What might have led to these results?

\section{Desiderata Revisited}
Answer based on the results. From a practical standpoints, were the desiderata fulfilled? Don't make mathematical proves in this chapter, the mathematical/logical aspects belong in "Modeling".

\section{Limitations in Experimental Setup}
What was not great about the evaluation (e.g. small sample size, only few combinations tested).

\section{Limitations of Metric}
Which technical limitations does the metric have? E.g. getting rid of the CE-Normalization (Payout Coefficient) for a Bayesian aproach or even choosing to break compatibility for better results. Be fair, but don't discredit my own work too much since this would pretty much destroy my thesis.

\section{Future Work}
Discuss what could be addressed in future work and what more needs to be researched?

\chapter{Conclusion}
Finish the entire thesis off in around 2, max 3 pages.
